\documentclass[11pt]{article}
\usepackage[T1]{fontenc}
\usepackage{textcomp}
\usepackage[margin=0.75in]{geometry}
\usepackage{amsmath,amssymb,amsthm}
\usepackage{array,booktabs}
\usepackage{hyperref}
\setlength{\parindent}{0pt}
\newtheorem{theorem}{Theorem}
\theoremstyle{definition}
\newtheorem{definition}{Definition}
\title{Flow Proof Artifact: prop-01-distributive-rectangle}
\author{Generated by \texttt{flow doc proof}}
\date{Flow Proof Book}
\begin{document}
\maketitle
If one of two straight lines is cut into segments, the rectangle by the whole pair equals the sum of rectangles by the uncut line and each segment.\\[0.5em]
\textbf{Source.} euclid — Elements, Book II, Proposition 1\\[0.5em]
\bigskip
\noindent{\large \textbf{Derived fact 1.} \textit{If one of two straight lines is cut into segments, the rectangle by the whole pair equals the sum of rectangles by the uncut line and each segment} \hfill the Euclidean plane \cdot Euclid Book II \cdot Proposition 1: the rectangle by two lines equals the sum over segments}\par
\smallskip\noindent\textit{Coordinate: the Euclidean plane \cdot Euclid Book II \cdot Proposition 1: the rectangle by two lines equals the sum over segments}\par
\smallskip\noindent\textit{Source: euclid — Elements, Book II, Proposition 1}\par
\smallskip\noindent\textit{Built on: proposition 34: complements of parallelograms about the diameter are equal, for Book I of the Elements in on the Euclidean plane}\par
\medskip
\noindent\textbf{Goal.} If one of two straight lines is cut into segments, the rectangle by the whole pair equals the sum of rectangles by the uncut line and each segment\par
\noindent$\displaystyle rect(A, B) = \text{sum rect}(A, \text{segments of B})$\par
\medskip
\noindent
\begin{tabular}{@{} >{\raggedright\arraybackslash}p{0.06\textwidth} >{\raggedright\arraybackslash}p{0.47\textwidth} >{\raggedleft\arraybackslash}p{0.41\textwidth} @{}}
\textbf{\#} & \textbf{Proof} & \textbf{Mathematics} \\
\midrule
\textbf{1.} & We prove that proposition 1: the rectangle by two lines equals the sum over segments for Euclid Book II the Euclidean plane. & \\[0.45em]
\textbf{2.} & Let A = uncut\_straight\_line. & \\[0.45em]
\textbf{3.} & Let B = straight\_line\_cut\_into\_segments. & \\[0.45em]
\textbf{4.} & Let parts = segments\_of\_B. & \\[0.45em]
\textbf{5.} & From step 2, step 3, and step 4, we can deduce that rect(A, B) equals sum rect(A, parts). & $\displaystyle rect(A, B) = \text{sum rect}(A, parts)$ \\[0.45em]
\textbf{6.} & We invoke the derived fact governing Book I of the Elements in the Euclidean plane: proposition 34: complements of parallelograms about the diameter are equal, for Book I of the Elements in on the Euclidean plane. & \\[0.45em]
\textbf{7.} & From step 6, this implies rect(A, B) equals sum rect(A, segments of B). Hence proven. & $\displaystyle rect(A, B) = \text{sum rect}(A, \text{segments of B})$ \\[0.45em]
\bottomrule
\end{tabular}
\label{thm:Geometry-Euclid Book II-proposition_1_the_rectangle_by_two_lines_equals_the_sum_over_segments}
\par\medskip\hrule
\end{document}
